\section{Scope, Source Base, and Method}

\subsection{What This Document Is Trying to Do}

The repository is not aiming to build a toy order-book bot.
The intended system is a research and trading platform for intraday market microstructure work across two primary domains:
\begin{itemize}
  \item crypto markets, with Bybit as the immediate live testing venue;
  \item Russian market instruments, especially those traded through T-Invest and the historical QScalp \texttt{.OrdLog.qsh} archive already stored in the repository.
\end{itemize}

The long-range goal is to move from informal discretionary pattern language toward:
\begin{itemize}
  \item explicit event definitions;
  \item reproducible logging and replay;
  \item measurable features;
  \item rule-based setup detectors;
  \item practical live research loops;
  \item and only later, if the data becomes good enough, ML overlays that filter or rank already-defined patterns.
\end{itemize}

This report therefore does \emph{three} jobs at once.
First, it preserves the original human trading logic.
Second, it organizes that logic into a coherent taxonomy.
Third, it translates each idea into a research agenda that can be implemented in code and tested against data.

\subsection{Source Inventory}

The knowledge base is a mixture of books and transcript-derived video materials.
The books provide higher-level structure and vocabulary; the transcripts contribute many of the practical details, caveats, and ``how it actually feels on the tape'' comments that are usually missing from books.

\subsubsection{Books}

\sourceentry{Professional Scalping}{Primary discretionary trading logic source for DOM, tape, clusters, densities, breakouts, bounces, POC, smart-money style framing, and robot behavior.}
\sourceentry{Prop Trading on MOEX}{Operational and infrastructure framing for prop-style intraday trading on the Russian market, including risk control, broker relationships, colocation, and platform requirements.}
\sourceentry{Risk Management in Trading}{Explicit risk-control framework for daily loss limits, per-trade risk, stop logic, leverage discipline, and the behavioral reasons traders break their own rules.}

\subsubsection{Transcript-Derived Video Materials}

The transcript series contributes concrete tactical material on:
\begin{itemize}
  \item densities;
  \item icebergs;
  \item tape and DOM configuration;
  \item robots;
  \item graph-plus-book interpretation;
  \item foundations for trades;
  \item bounce trades;
  \item breakout trades;
  \item trade potential, exits, and adds;
  \item decision logic;
  \item currency and futures specifics;
  \item limit-up / limit-down regimes;
  \item springs and stretched moves;
  \item news-driven behavior;
  \item and self-analysis of trading.
\end{itemize}

Two transcript files in the repository are effectively empty or near-empty:
\begin{itemize}
  \item the volatility-harvesting lesson;
  \item the psychology lesson.
\end{itemize}
The relevant ideas from those topics were therefore reconstructed mostly from adjacent materials, not from fully populated transcripts.

\subsection{Translation Philosophy}

The source material is discretionary and heavily contextual.
That matters.
Many ideas are described in a way that makes perfect sense to an experienced human staring at a fast DOM, but not to a research pipeline.
Phrases such as ``clean level'', ``good participant'', ``price is being held'', ``the book is empty'', or ``today the instrument is moving properly'' are meaningful, but they are not directly measurable.

The translation philosophy used in this report is:
\begin{enumerate}
  \item preserve the original meaning instead of prematurely forcing false precision;
  \item separate what is explicit from what is fuzzy;
  \item identify observables, even if the threshold values remain open;
  \item map discretionary logic into candidate state machines and feature groups;
  \item state clearly when the available data is insufficient to test the full human concept.
\end{enumerate}

\subsection{Important Reliability Caveats}

This report should \emph{not} be read as proof that the described setups are profitable.
It is a formalization of practitioner knowledge, not a performance report.

Several reliability issues must stay visible:
\begin{itemize}
  \item The books and videos represent a specific trading style, with a strong DOM-and-tape bias.
  \item Some examples are tied to particular MOEX eras, liquidity profiles, and terminal workflows.
  \item The historical Russian archive in the repository is order-book only, old, and often sparse or low quality.
  \item A meaningful portion of the original edge may depend on data not present in that archive, especially trades, execution response, and hidden-liquidity inference.
  \item Human discretion often uses cross-context cues that are easy to mention and hard to encode.
\end{itemize}

The right use of this document is therefore:
\begin{itemize}
  \item as a hypothesis map;
  \item as a design document for logging and replay;
  \item as a guide for detector implementation;
  \item and as a record of what still requires live validation.
\end{itemize}
